% AI collaboration: orchestrator and roles, and a selected prompt with its
% result and evaluation.
% The prompt excerpt quotes the real stage 6a task card (ai-log 2026-11-02).

\begin{frame}
  \frametitle{Każdą pracę agenta sprawdzał recenzent i orkiestrator}
  \begin{itemize}
    \item Sesja główna Claude Code jest orkiestratorem i deleguje zadania
      kartami z celem, wycinkiem wymagań z referencjami, listą plików
      i kryteriami akceptacji.
    \item Pracę coder sprawdza code-reviewer, a pracę latex-assistant
      presentation-reviewer. Definicje ról recenzentów dają tylko narzędzia
      odczytu, ale przy tej talii ograniczenie wynikało z polecenia.
    \item Kompilację, kontrolę czcionek i niezależne obliczenia kontrolne
      uruchamia sam orkiestrator i przekazuje ich wynik recenzentom.
    \item AI proponuje decyzje i tezy slajdów, a autor zatwierdza je
      w punktach kontrolnych.
  \end{itemize}
\end{frame}

\begin{frame}[fragile]
  \frametitle{Kryterium w liczbach nie zostawia agentowi pola do interpretacji}
\begin{Verbatim}[commandchars=\\\{\},fontsize=\footnotesize]
[...]
Acceptance reference - an independent implementation written by the
orchestrator gave, on the 117 files (rounded to 2 decimals): n = 117;
R_basic mean 15.65, median \textminus{}0.30; \ensuremath{\Delta}R mean \textminus{}4.16 pp, median 0.00 pp;
helped 49, hurt 43, same 25; buy-and-hold mean 43.33, median 14.95.
Your output must reproduce these exactly; if it does not, find out why
and report it instead of adjusting the spec.
[...]
\end{Verbatim}
  \vskip1mm
  Wynik: agent coder w obliczeniach eksperymentu odtworzył wszystkie wartości dokładnie.\\
  Ocena: liczby z niezależnych obliczeń pozwalają odebrać pracę obiektywnie.
\end{frame}
